\subsection{A quantitative estimate on curves}
\label{sec:ch28-curve-stability}

For measures carried by one finite curve skeleton, the variational problem
is a weighted graph Dirichlet problem.  Stability becomes an explicit
inequality.  We keep the total modulus length in the stated constants,
while also recording the sharper resistance-length bound.

Throughout this subsection we retain \eqref{eq:ch28-standing}:
\(X\) is a smooth connected projective curve over the original discretely
valued field, \(\mathcal X\) is a regular semistable model with skeleton
\(\Gamma\), \(r:X^{\An}\to\Gamma\) is its retraction, and \(\omega\)
is a semipositive \emph{real} closed model form of ample generic class
determined on \(\mathcal X\).  Its measure \(\mu_0=\MA_\omega(0)\)
has mass \(V=(\{\omega\})>0\).  No split-node or geometric-connectedness
assumption is added.  Give each edge its modulus length \(\ell_e\) and
its residue weight \(c_e=[\kappa(z_e):\widetilde K]\geq1\), retained
on subdivisions, and write
\begin{equation}
\label{eq:ch28-dirichlet-form}
 D_\Gamma(f):=\sum_e c_e\int_e|f'|^2\,\rd s\ \in[0,+\infty],
 \qquad L(\Gamma):=\sum_e\ell_e,
 \qquad L_R(\Gamma):=\sum_e\frac{\ell_e}{c_e}\leq L(\Gamma).
\end{equation}
This defines the Dirichlet energy for continuous functions that are
absolutely continuous on each edge, with value \(+\infty\) if the
weighted derivative integral diverges.  The energy is weighted but the
slope coordinate is unchanged.  Both lengths and the energy are zero
when the graph is a point.

\begin{lemma}
\label{lem:ch28-graph-sobolev}
Let \(\Gamma\) be a compact connected finite metric graph with positive
edge weights \(c_e\geq1\), including the unweighted case \(c_e=1\).
If \(f\) is continuous and absolutely continuous on each edge, then
\begin{equation}
\label{eq:ch28-graph-sobolev}
 \operatorname{osc}_\Gamma f
 \leq L_R(\Gamma)^{1/2}D_\Gamma(f)^{1/2}
 \leq L(\Gamma)^{1/2}D_\Gamma(f)^{1/2}.
\end{equation}
In particular this holds for the originally stated piecewise \(C^1\)
functions.
\end{lemma}

\begin{proof}
The conclusion is immediate if the graph is a point or the energy is
infinite.  Otherwise join \(x,y\) by an injective path \(\gamma\),
subdividing at its endpoints.  Such a path exists by connectedness and
traverses no edge segment twice.  The fundamental theorem of calculus and
weighted Cauchy--Schwarz give
\[
 |f(x)-f(y)|
 \leq\int_\gamma|f'|\,\rd s
 \leq\left(\int_\gamma\frac{\rd s}{c_e}\right)^{1/2}
       \left(\int_\gamma c_e|f'|^2\,\rd s\right)^{1/2}
 \leq L_R(\Gamma)^{1/2}D_\Gamma(f)^{1/2}.
\]
Taking the supremum and using \(L_R(\Gamma)\leq L(\Gamma)\) proves
\eqref{eq:ch28-graph-sobolev}.
\end{proof}

\begin{theorem}
\label{thm:ch28-curve-stability}
\index{stability!with respect to the measure}
Let \(\mu_1,\mu_2\) be positive measures of mass \(V\) carried by
\(\Gamma\), and let \(u_1,u_2\) be the continuous BFJ solutions of
\(\MA_\omega(u_a)=\mu_a\) constructed in
\cref{lem:ch26-weighted-graph-real-bfj}, normalized by
\(\sup_{X^{\An}}u_a=0\).  Put \(u=u_1-u_2\), and write
\(\lVert\mu_1-\mu_2\rVert\) for total variation on \(\Gamma\).  Then
\begin{align}
\label{eq:ch28-curve-stability-dirichlet}
 D_\Gamma(u)&\leq L(\Gamma)\lVert\mu_1-\mu_2\rVert^2,\\
\label{eq:ch28-curve-stability-sup}
 \lVert u_1-u_2\rVert_{C^0(X^{\An})}
   &\leq L(\Gamma)\lVert\mu_1-\mu_2\rVert,\\
\label{eq:ch28-curve-stability-dual}
 \bigl|E^\vee_\omega(\mu_1)-E^\vee_\omega(\mu_2)\bigr|
   &\leq\tfrac12L(\Gamma)\lVert\mu_1-\mu_2\rVert^2
      +L(\Gamma)^{1/2}\lVert\mu_1-\mu_2\rVert
          \min_a D_\Gamma(u_a)^{1/2}.
\end{align}
The last bound is written without inverse powers of \(L(\Gamma)\),
so it also makes sense for a point skeleton.
Each occurrence of \(L(\Gamma)\) in the three bounds may be replaced
by \(L_R(\Gamma)\).  In particular \(\mu\mapsto u_\mu\) is Lipschitz
from total variation to the uniform norm with the stated constant, and
\(E^\vee_\omega\) is locally Lipschitz on measures carried by \(\Gamma\).
\end{theorem}

\begin{proof}
\emph{The finite-energy identity.}
The construction in \cref{lem:ch26-weighted-graph-real-bfj} applies to
the real form \(\omega\) and provides the asserted BFJ solutions over
the original field.  It also proves edgewise absolute continuity,
finite weighted energy, and integration by parts for their difference.
Since \(\Delta_\Gamma u=\mu_1-\mu_2\),
\begin{equation}
\label{eq:ch28-stability-dirichlet-identity}
 D_\Gamma(u)=-\int_\Gamma u\,\rd\Delta_\Gamma u
            =-\int_\Gamma u\,\rd(\mu_1-\mu_2).
\end{equation}
Continuity alone would not have justified the derivative integral; the
strong derivative convergence in that lemma does.

\emph{The estimate.}
The signed measure has mass zero.  Subtract the midpoint of the range of
\(u\) before bounding its integral, and apply
\cref{lem:ch28-graph-sobolev}:
\[
 D_\Gamma(u)
 \leq\tfrac12\operatorname{osc}_\Gamma(u)\lVert\mu_1-\mu_2\rVert
 \leq\tfrac12 L_R(\Gamma)^{1/2}D_\Gamma(u)^{1/2}
                       \lVert\mu_1-\mu_2\rVert.
\]
If \(D_\Gamma(u)=0\) there is nothing to divide by.  Otherwise divide
by its square root.  In both cases,
\[
 D_\Gamma(u)^{1/2}
 \leq\tfrac12L_R(\Gamma)^{1/2}\lVert\mu_1-\mu_2\rVert,
 \qquad
 \operatorname{osc}_\Gamma u
 \leq\tfrac12L_R(\Gamma)\lVert\mu_1-\mu_2\rVert.
\]
This proves the first asserted bound, with room to spare.

\emph{The uniform norm on the analytic curve.}
Both solutions factor through \(r\), so their difference has the same
uniform norm on \(X^{\An}\) and on \(\Gamma\).  Each restriction has
maximum zero.  At a maximum of \(u_2\), the difference is nonpositive;
at a maximum of \(u_1\), it is nonnegative.  Thus
\(\lVert u_1-u_2\rVert_\infty\leq\operatorname{osc}_\Gamma(u_1-u_2)\).
The preceding estimate proves the second bound.  No factorization of
an arbitrary higher-dimensional solution is being asserted.

\emph{The dual energies.}
The independent calculation in \cref{prop:ch28-curve-dual-energy} below
uses \cref{lem:ch26-weighted-graph-real-bfj,prop:ch26-graph-energy}, not
the present theorem, and gives \(E^\vee_\omega(\mu_a)=D_\Gamma(u_a)/2\).
Put \(A_a=D_\Gamma(u_a)^{1/2}\) and \(B=D_\Gamma(u_1-u_2)^{1/2}\).
The weighted derivative norm is a seminorm, so
\[
 |A_1-A_2|\leq B,\qquad A_1+A_2\leq2\min(A_1,A_2)+B.
\]
Consequently
\[
 \bigl|E^\vee_\omega(\mu_1)-E^\vee_\omega(\mu_2)\bigr|
 \leq\tfrac12B^2+B\min(A_1,A_2).
\]
Use the weaker bound
\(B\leq L_R(\Gamma)^{1/2}\lVert\mu_1-\mu_2\rVert\) already proved.
This is the third estimate with \(L_R\), hence also with \(L\).
The same estimates bound \(A_a\) on every total-variation bounded
neighborhood of a fixed measure, proving the local Lipschitz assertion.
For a point skeleton the sole measure of mass \(V\) is \(\mu_0\),
the normalized solution is zero, and all three estimates read \(0\leq0\).
\end{proof}

\begin{remark}[What the constant measures, and what replaces it]
\label{rem:ch28-curve-constant}
The constant \(L(\Gamma)\) is the total analytic modulus length.
The sharper \(L_R(\Gamma)\) is the total resistance length
\(\sum_e\ell_e/c_e\).  Subdividing edges preserves both totals, because
the sublengths sum to the old length and each piece retains \(c_e\).
General domination may also attach branches, so it need not preserve
either total.  For a function pulled back from the old skeleton, its
derivative on an attached branch is zero and its weighted energy is
unchanged.  One may retain the estimate on the old common support
skeleton or use the possibly larger constant on the new one.

There is no unconditional rule that finite scalar extension multiplies
every analytic length by the ramification index.  With the absolute value
extended from \(K\), the modulus \(-\log|\alpha|\) of a given lifted
node chart \(xy=\alpha\) is unchanged; splitting, new components and
residue degrees must still be accounted for.  If instead the discrete
valuation is renormalized so that the new uniformizer has value one,
all logarithmic lengths and potentials are being rescaled.  That is a
different normalization, not the fixed-absolute-value scalar extension.
Apply the estimates to the actual resulting weighted graph and the
consistent measure conventions of \cref{thm:base-change-ma}.

In dimension \(n\geq2\) the scalar Dirichlet coefficient is replaced by
several mixed coefficients.  A bound for one of them does not control
oscillation, for the reason in \cref{rem:ch28-uniqueness-gap}.
The qualitative \cref{prop:fixed-skeleton-stability} and the
measure estimate \cref{thm:finite-energy-measure-stability} remain the
separately stated higher-dimensional results.
\end{remark}
